\scanpage{16}
profinite, to a certain extent!) use the geometric language of the fundamental groupoids of $\mathbb P^1(\mathbb C)-\{0,1,\infty\}$. We \deleted{\unsure{shall be able to|shall set}} shall therefore denote by $\Pi=\Pi_{0,3}$ the discrete group
\[
\Pi=\Pi_{\unsure{1|*}}(\overbrace{\mathbb P^1(\mathbb C)\setminus\{0,1,\infty\}}^{U=U_{0,3}},P)
\]
where $\mathbb P^1(\mathbb C)$ is endowed with the unique \inserted{oriented} Riemannian \inserted{sphere} structure for which the rotation \inserted{$\dot\rho$} through an angle of $+\frac{2\pi}{3}$ about the axis $P,\bar P$ \deleted{\unsure{which}} induces on $(0,1,\infty)$ the cyclic permutation $0\longrightarrow1\longrightarrow\infty\longrightarrow0$ This corresponds to the usual orientation of the sphere), and where the complex automorphisms \inserted{$\dot\sigma_0,\dot\sigma_1,\dot\sigma_\infty$} of $\mathbb P^1_{\mathbb C}$ which \deleted{\unsure{exchange}} \inserted{\deleted{\unsure{fix}}} \unsure{induce|two $p^t$} \deleted{\unsure{resp|on}.} on $(0,1,\infty)$ the transpositions fixing resp. $0,1,\infty$, are \inserted{resp.} the \deleted{\unsure{rotations}} \inserted{symmetries} about the axes \deleted{$\unsure{P,\bar P|1,\infty|0,1,\infty}$} $(0,2)$, $(1,-1)$ and $(\infty,\frac12)$.
\margin{On the Riemann sphere we have
\[
\begin{aligned}
\sigma_0(z)&=\frac{-\unsure{z|1}}{1-z}=\frac{z}{z-1}\\
\sigma_1(z)&=\frac1z\\
\sigma_\infty(z)&=1-z\\
\rho(z)&=\frac1{1-z}\\
\rho^2(z)&=1-\frac1z
\end{aligned}
\]}
As for complex conjugation, it is the symmetry with respect to \unsure{the|a} plane $(0,1,\infty)$. We consider $\unsure{\mathfrak S_3|\mathcal S_3}$ as acting on $(0,1,\infty)$ \unsure{in|in} this order, and $\unsure{\mathfrak S_3|\mathcal S_3}\times\mathbb Z/2$ as acting on the sphere; via the $\dot\sigma_0,\dot\sigma_1,\dot\sigma_\infty\in\unsure{\mathfrak S_3|\mathcal S_3}$ and $\dot\tau\in\mathbb Z/2$. We consider the extension \inserted{$\deleted{\unsure{\mathcal E_0}}$} of $\unsure{\mathfrak S_3|\mathcal S_3}\times\mathbb Z/2$ by $\Pi$, denoted \inserted{$\mathcal E_{0,3}$ or} $\mathcal E$, defined by
\[
\mathcal E=\Pi_1\left(\underset{\substack{\Vert\\U_{0,3}}}{U_{\deleted{\unsure{0,3}}}},\underbrace{\unsure{\mathfrak S_3|\mathcal S_3}\times\mathbb Z/2}_{\Gamma};P\right)
\]
\mspage{15}
